\section{Related work}
\label{sec:related}
\para{Open-ended Environments for Decision-making Agents.} 
There are many environments developed with the goal of open-ended agent learning. Prior works include maze-style worlds \cite{team2021openended,wang2019paired, juliani2019obstacletower}, purely text-based game~\cite{heinrich2020nethack}, grid worlds ~\cite{boisvert2019babyai,cao2020babyai}, browser/GUI-based environments ~\cite{shi2017worldofbits,toyama2021androidenv}, and indoor simulators for robotics ~\cite{deepmind2020playroom,shen2020igibson,srivastava2021behavior,fan2021secant,shridhar2020alfred,savva2019habitat,puig2018virtualhome}. 
Minecraft offers an exciting alternative for open-ended agent learning. It is a 3D visual world with procedurally generated landscapes and extremely flexible game mechanics that support an enormous variety of activities.
Prior methods in open-ended agent learning \Rebuttal{\mbox{\cite{  ecoffet2019goexplore, huizinga2022evolving, wang2020enhanced,  kanitscheider2021multitask, dennis2020paired}} do not make use of external knowledge, but our approach leverages internet-scale database to learn open-vocabulary reward models, thanks to Minecraft's abundance of gameplay data online}.  

\vspace{-0.05in}
\para{Minecraft for AI Research.} The Malmo platform \cite{johnson16malmo} is the first comprehensive release of a Gym-style agent API \cite{brockman2016openai} for Minecraft. Based on Malmo, MineRL \cite{guss2019minerl} provides a codebase and human play trajectories for the annual Diamond Challenge at NeurIPS \cite{guss2019minerl2,guss2021minerl,kanervisto2022minerl}. \minedojo's simulator builds upon the pioneering work of MineRL, but greatly expands the API and benchmarking task suite. Other Minecraft benchmarks exist with different focuses. For example, CraftAssist~\cite{gray2019craftassist} and IGLU~\cite{kiseleva2021neurips} study agents with interactive dialogues.  BASALT~\cite{shah2021basalt} applies human evaluation to 4 open-ended tasks. EvoCraft~\cite{grbic21evocraft} is designed for structure building, and Crafter~\cite{hafner2021crafter} optimizes for fast experimentation. Unlike prior works, \minedojo's core mission is to facilitate the development of generally capable embodied agents using internet-scale knowledge. \Rebuttal{We include a feature comparison table of different Minecraft platforms for AI research in Table~\ref{supp:table:framework_comparison}}.

\vspace{-0.05in}
\para{Internet-scale Multimodal Knowledge Bases.} 
Big dataset such as Common Crawl~\cite{commoncrawl}, the Pile \cite{gao2020pile}, LAION \cite{schuhmann2021laion}, YouTube-8M \cite{haija2016youtube8m} and HowTo100M~\cite{miech2019howto100m} have been fueling the success of large pre-trained language models~\cite{devlin2018bert,radford2019gpt2, brown2020gpt3} and multimodal models~\cite{sun2019learning, DBLP:conf/nips/AlayracRSARFSDZ20, DBLP:conf/cvpr/MiechASLSZ20, zhu2020actbert, DBLP:conf/aaai/AmraniBRB21, akbari2021vatt, xu2021videoclip}. While generally useful for learning representations, these datasets are not specifically targeted at embodied agents. To provide agent-centric training data, RoboNet~\cite{dasari19robonet} collects video frames from 7 robot platforms, and Ego4D~\cite{grauman2021ego4d} recruits volunteers to record egocentric videos of household activities. In comparison, \minedojo's knowledge base is constructed without human curation efforts, much larger in volume, more diverse in data modalities, and comprehensively covers all aspects of the Minecraft environment. 

\vspace{-0.05in}
\para{Embodied Agents with Large-scale Pre-training.} 
Inspired by the success in NLP, embodied agent research\Rebuttal{~\mbox{\cite{duan2022survey, batra2020rearrangement, ravichandar2020recent, collins2021review}}} has seen a surge in adoption of the large-scale pre-training paradigm. The recent advances can be roughly divided into 4 categories. 
1) \textbf{Novel agent architecture}: Decision Transformer \cite{ranzato2021decisiontransformer,janner2021onebigsequence,zheng2022onlinedt} applies the powerful self-attention models to sequential decision making. \Rebuttal{GATO~\mbox{\cite{reed2022generalist}} and Unified-IO~\mbox{\cite{lu2022unifiedio}} learn a single model to solve various decision-making tasks with different control interfaces.} VIMA~\cite{jiang2022vima} unifies a wide range of robot manipulation tasks with multimodal prompting. 
2)  \textbf{Pre-training for better representations}: R3M~\cite{nair2022r3m} trains a general-purpose visual encoder for robot perception on Ego4D videos~\cite{grauman2021ego4d}. CLIPort~\cite{shridhar2021cliport} leverages the pre-trained CLIP model \cite{radford2021clip} to enable free-form language instructions for robot manipulation. 
3) \textbf{Pre-training for better policies}: AlphaStar \cite{vinyals2019alphastar} achieves champion-level performance on StarCraft by imitating from numerous human demos. SayCan \cite{ahn2022saycan} leverages large language models (LMs) to ground value functions in the physical world. \cite{li2022pretrained} and \cite{reid2022wikipedia} directly reuse pre-trained LMs as policy backbone. \Rebuttal{VPT~\mbox{\cite{openai2022vpt}} is a concurrent work that learns an inverse dynamics model from human contractors to pseudo-label YouTube videos for behavior cloning. VPT is complementary to our approach, and can be finetuned to solve language-conditioned open-ended tasks with our learned reward model. } 
4) \textbf{Data-driven reward functions}: Concept2Robot~\cite{shao2021concept2robot} and DVD~\cite{nair2021dvd} learn a binary classifier to score behaviors from in-the-wild videos~\cite{goyal2017something}. LOReL~\cite{nair2021lorel} crowd-sources humans labels to train language-conditioned reward function for offline RL. AVID \cite{smith2019avid} and XIRL \cite{zakka2021xirl} extract reward signals via cycle consistency. 
\minedojo's task benchmark and internet knowledge base are generally useful for developing new algorithms in all the above categories. In Sec. \ref{sec:method}, we also propose an open-vocabulary, multi-task reward model using \minedojo YouTube videos.

